\subsection{The Least-Squares Optimality of Webster}

\begin{theorem}[Webster Least-Squares Optimality]
  Among all apportionments satisfying $\sum_i s_i = H$ and $s_i \geq 1$,
  Webster minimises the sum of squared deviations from exact quotas:
  \[
    \text{Webster} = \arg\min_{\{s_i\}} \sum_{i=1}^{n} (s_i - q_i)^2,
  \]
  where $q_i = p_i \times H / N$ is state $i$'s exact proportional quota.
\end{theorem}

\begin{proof}[Proof sketch]
  Consider a transfer of one seat from state $j$ to state $i$
  (where $s_j \geq 2$ and $s_i \geq 1$).
  The change in $\sum (s_k - q_k)^2$ is:
  \[
    \Delta = (s_i + 1 - q_i)^2 + (s_j - 1 - q_j)^2
           - (s_i - q_i)^2 - (s_j - q_j)^2.
  \]
  The transfer reduces the sum iff $\Delta < 0$, i.e.,
  iff $q_i - s_i > q_j - s_j + 1 - 1 = q_j - s_j$,
  i.e., iff state $i$ has a larger fractional excess from its
  current allocation than state $j$.
  At the Webster optimum, no such improving transfer exists,
  meaning each state's residual $q_i - s_i$ is within $[-0.5, 0.5]$,
  which is precisely the condition for $s_i = \lfloor q_i + 0.5 \rfloor$.
  (More precisely, the proof uses the fact that integer rounding at 0.5
  is the unique solution to the least-squares nearest-integer problem.)
  See \citet{balinski1982fair}, ch.~3, for the full proof.
\end{proof}

This optimality property makes Webster the natural ``proportionality
benchmark'' in academic apportionment analysis:
it is the divisor method that gets closest to strict proportionality
in the mean-squared sense.

\subsection{Neutrality Between Small and Large States}

Webster is \emph{neutral} between small and large states in the following
precise sense: the arithmetic mean rounding threshold $s + 0.5$ is the
symmetric midpoint between $s$ (no extra seat) and $s+1$ (extra seat).
States with exact quotas above the midpoint receive more; those below
receive fewer.
This symmetry means that, in expectation over uniformly distributed
exact quotas, Webster neither systematically overrepresents nor
underrepresents any size class of states.

By contrast:
\begin{itemize}
  \item \textbf{Adams} (ceiling rounding): every state with a fractional
    quota above $s$ receives $s+1$ seats, giving small states with
    quotas barely above an integer the same benefit as large states
    with quotas nearly at the next integer.
    This systematically overrepresents small states.
  \item \textbf{Huntington-Hill}: threshold $\sqrt{s(s+1)} < s+0.5$,
    so states need a \emph{smaller} fractional excess to round up
    under HH than under Webster.
    The excess benefiting small states is small but nonzero.
  \item \textbf{Jefferson} (floor rounding): every fractional part is
    lost, systematically underrepresenting small states
    (which have large fractional excesses relative to their seat counts).
\end{itemize}

The ordering of bias (from most small-state-favoring to most
large-state-favoring) is:
\[
\text{Adams} \succ \text{HH} \succ \text{Webster} \succ \text{Jefferson}.
\]
Webster is the unique divisor method with zero systematic bias
in the arithmetic-mean sense.

\subsection{Paradox Immunity}

Webster, as a divisor method, is immune to the Alabama paradox,
the population paradox, and the new-states paradox.
The proof is identical in structure to the proof for Huntington-Hill
(J.1, Theorem~2): the priority queue monotonically allocates seats
based on population, so increasing $H$ or increasing a state's
population can only add seats, never remove them.
The full paradox immunity proof for all divisor methods is given
in J.5.
